% TOP_PROBLEM: {"id": 20000914, "title": "A high-probability lower bound for random permutation discrepancy", "category_id": 2, "category": {"id": 2, "name": "combinatorics", "display_name": "Combinatorics", "description": "Counting problems, graph theory, discrete structures.", "slug": "combinatorics", "order_index": 2, "created_at": "2026-07-31T15:26:25.671Z"}, "status": "partially_solved", "problem_number": "AIM-COMBINATORICS-0039", "set_id": 14, "statusQualification": "Uses all permutation intervals; the prefix convention changes discrepancy by at most a factor of two."}
% FIRST_POSTED: 2026-10-01
% ENTRY_KIND: statement_only
% UnsolvedMath ID 20000914; source catalogue code AIM-COMBINATORICS-0039
% !TeX program = lualatex
% Definition and sources only; this file is not an LLM solution attempt.
\documentclass[11pt,a4paper]{article}
\usepackage[margin=25mm]{geometry}
\usepackage{fontspec}
\setmainfont{lmroman10-regular.otf}[BoldFont=lmroman10-bold.otf,
 ItalicFont=lmroman10-italic.otf,BoldItalicFont=lmroman10-bolditalic.otf]
\usepackage{amsmath,amssymb,mathtools}
\usepackage{unicode-math}
\setmathfont{latinmodern-math.otf}
\AtBeginDocument{\let\setminus\smallsetminus}
\usepackage{xurl,hyperref}
\hypersetup{colorlinks=true,urlcolor=blue}
\setcounter{secnumdepth}{0}
\setlength{\parindent}{0pt}
\setlength{\parskip}{5pt}
\setlength{\emergencystretch}{3em}
\begin{document}
\section*{A high-probability lower bound for random permutation discrepancy}\label{20000914}
{\small UnsolvedMath ID 20000914 · AIM-COMBINATORICS-0039 · Combinatorics · AIM Workshop Problem Lists}
\subsection{Definitions and mathematical statement}
Write $[n]=\{1,\ldots,n\}$. A permutation $\pi$ of $[n]$ gives the
interval sets $\{\pi(a),\pi(a+1),\ldots,\pi(b)\}$ for
$1\le a\le b\le n$. Given $k$ permutations
$\pi=(\pi_1,\ldots,\pi_k)$, take the union of these $k$
interval systems and define
\[
 D(\pi)=\min_{\varepsilon:[n]\to\{-1,1\}}
 \max_{\substack{1\le i\le k\\1\le a\le b\le n}}
 \left|\sum_{j=a}^b\varepsilon(\pi_i(j))\right|.
\]
There is one sign for each ground-set element, used consistently in all
permutations. Repeated permutations and repeated interval sets cause no
change to the definition.

The first question is to determine the worst-case discrepancy
$D(n,k)=\max_{\pi}D(\pi)$, with sharp dependence
on both the ground-set size $n$ and the number $k$ of permutations.
The second asks for the typical value of $D(\pi)$ when
$\pi_1,\ldots,\pi_k$ are independent uniform random permutations,
including asymptotic bounds holding with probability tending to one as
$n\to\infty$ for the specified regime of $k=k(n)$.
These are distinct extremal and probabilistic questions; a bad fixed
family need not describe a random family. Intervals are the convention
here. If one instead uses only prefixes, its discrepancy lies between
one half and one times the interval discrepancy, so the asymptotic
order questions are unchanged.
\subsection{Short English statement}
Determine the worst-case and typical discrepancy of the interval systems generated by several permutations.
\subsection{Sources}
\begin{itemize}
\item[S1] ULAM.AI and the UnsolvedMath Contributors, supplied problems.json, record 20000914 (AIM-COMBINATORICS-0039), AIM Workshop Problem Lists. Catalogue identity and source transcription; CC BY 4.0.
\url{https://huggingface.co/datasets/ulamai/UnsolvedMath}.
\item[S2] American Institute of Mathematics, Hereditary discrepancy and factorization norms, item 1.3. Original workshop question underlying AIM-COMBINATORICS-0039.
\url{http://aimpl.org/hereddiscrep/1/}.
\end{itemize}
\end{document}
